% Plantilla de analysis.pdf (Quiz 03). 
% Los comentarios % bajo cada apartado son pistas: no salen en el PDF.
% Limite de la entrega: 4 paginas.
% Compilar: pdflatex analysis_plantilla.tex
\documentclass[11pt,a4paper]{article}
\usepackage[utf8]{inputenc}
\usepackage[T1]{fontenc}
\usepackage[spanish,es-noshorthands]{babel}
\usepackage{amsmath,amssymb}
\usepackage{booktabs,array}
\usepackage{enumitem}
\usepackage{titlesec}
\usepackage[a4paper,top=1.4cm,bottom=1.5cm,left=1.7cm,right=1.7cm]{geometry}

\setlength{\parindent}{0pt}
\setlength{\parskip}{2pt}
\setlength{\emergencystretch}{2em}
\setlist{nosep,leftmargin=1.3em}
\titleformat{\section}{\large\bfseries}{}{0pt}{}
\titlespacing*{\section}{0pt}{9pt}{2pt}
\newcolumntype{L}[1]{>{\raggedright\arraybackslash}p{#1}}
\newcommand{\phys}{\operatorname{phys}}
\newcommand{\sub}[2]{\medskip\par\noindent\textbf{#1) #2.}\ }

\begin{document}

\begin{center}
{\Large\bfseries Estructuras de Datos --- Quiz 03}\\[1pt]
{\large Dynamic Arrays and Amortized Analysis} \quad$\cdot$\quad analysis.pdf \quad$\cdot$\quad Semestre 2026-2
\end{center}
\vspace{-6pt}
\textbf{Nombres:} Kevin Javier Gonzalez, Iván Felipe Maluche Suarez, Angelica Pascagaza Vega \hfill \textbf{Documento:} 1093432420 \hfill \textbf{Grupo:} 3

\section*{Notación y modelo de costo}

%=====================================================================
\section*{P1 --- Arreglo dinámico desde cero}

\sub{a}{Implementación}
% Referencia a src/DynamicArray.java. Representacion: data, size, capacity.
En el archivo de DynamicArray.Java se encuentra la implementacion del arreglo dinamico.

\sub{b}{Invariantes}

\begin{enumerate}[label=\textbf{I\arabic*.}]
\item $$0 \leq size \leq capacity$$
\item $$ capacity \geq 1$$
\item $$ index \in [0, size), data[size] == null $$
\end{enumerate}
% Justificar la preservacion en append / removeLast / resize.


\sub{c}{Complejidad worst-case}\par\smallskip
\renewcommand{\arraystretch}{2}
\begin{tabular}{@{}L{0.16\textwidth}L{0.2\textwidth}L{0.55\textwidth}@{}}
\toprule
Operación & Worst-case & Justificación\\
\midrule
\texttt{get}       & $\Theta(1)$  & El metodo solamente comprueba si el index es valido y devuelve el dato, debido a que el array por defecto ya lo hace sin necesidad de recorrer nada cuesta O(1) \\ 
\texttt{set}        & $\Theta(1)$ & Similar a get, es una busqueda de costo 1 mas en este caso se asigna al valor un dato manteniendo constante la cota\\
\texttt{append}     & O(n) & Debido al resize en el peor caso es necesario construir un nuevo array y dato por dato del anterior asignarlo en el mismo.\\
\texttt{removeLast} & $\Theta(1)$ & Tal como el caso de get y set consiste en guardar el ultimo valor, asginarle un nulo y disminuir size, siendo operaciones que me dan una cota constante.\\
\bottomrule
\end{tabular}
\renewcommand{\arraystretch}{1}


\sub{d}{Worst-case vs.\ amortizado}

Al realizar el append siempre que $size$ < a $capacity$ el coste de operacion sera "barata", lo cual permite que cuando pase el caso en que ($size = capacity$) el costo de resize se amoritugara en el total.

%=====================================================================
\section*{P2 --- Snake sobre \texttt{DynamicArray} con desplazamiento}

\sub{a}{Costo de \texttt{removeTail}}

Al llamar removeTail en esa implementacion es necesario eliminar la cola y luego reasignar los valores que quedan una posicion hacia atras, esto implica un costo O(n) (con n siendo en este caso el size) al tener que recorrer el array .

\sub{b}{Costo total de $m$ movimientos}

Al ser de longitud k habria un costo de $O(k)$ lo cual al realizarse m veces nos daria un costo total de $O(km)$.

\sub{c}{Validez de la afirmación} % por hacer
Esto es falso, $move()$ no es amortizado porque en cada operacion de $move()$ tengo que hacer n (size) actualizaciones de posicion de la serpiente lo que hace que en promedio sea $\frac{O(n \cdot n)}{n} \rightarrow O(n)$ veces el costo amortizado, a diferencia del caso de append donde gracias a la distribucion entre operaciones caras y baratas el coste se amortigua.


\sub{d}{Bottleneck} % por hacer
% Que fija la representacion y por que obliga a mover k elementos.

El bottleneck en esta implementacion consistiria en que debido a que en cada movimiento que la serpiente no consuma una manzana (lo cual pasara frecuentemente) sera necesario ejecutar un RemoveTail, el cual como se explico en el punto anterior es O(n).
Esto culminara en que esta operacion basica y siempre necesaria para todo el loop jugable tenga un costo considerable y pueda generar problemas de estabilidad al ejecutarse en casi todo momento.

%=====================================================================
\section*{P3 --- Rediseño: representación interna}

\subsection*{a) Variables de estado}

\begin{itemize}[leftmargin=1.2cm]
    \item \texttt{data}: arreglo nativo \texttt{Object[]} de tamaño \texttt{capacity} donde se almacenan las posiciones.
    \item \texttt{tail}: índice físico donde está la \textbf{cola} (índice lógico $0$).
    \item \texttt{size}: número de elementos vivos (longitud lógica de la serpiente).
    \item \texttt{capacity}: tamaño físico del arreglo \texttt{data}.
\end{itemize}

La cabeza está en $(\texttt{tail} + \texttt{size} - 1) \bmod \texttt{capacity}$. Solo se necesitan estas cuatro variables (constante).

\subsection*{b) Mapeo lógico}

\[
\text{phys}(i) = (\texttt{tail} + i) \bmod \texttt{capacity}, \qquad 0 \le i < \texttt{size}.
\]

Propiedades:
\begin{enumerate}[leftmargin=1.2cm]
    \item \textbf{Rango:} como $\texttt{tail} \in [0, C)$ y $i \in [0, \texttt{size}) \subset [0, C)$, se cumple $0 \le \text{phys}(i) < C$.
    \item \textbf{Inyectividad:} si $\text{phys}(i) = \text{phys}(j)$ con $0 \le i, j < \texttt{size} \le C$, entonces $i \equiv j \pmod{C}$, y como $|i - j| < C$, se tiene $i = j$.
    \item \textbf{Celda libre:} si $\texttt{size} < C$, $\text{phys}(\texttt{size}) = (\texttt{tail} + \texttt{size}) \bmod C$ no coincide con ningún $\text{phys}(i)$ con $i < \texttt{size}$, así que esa celda está libre para el próximo \texttt{addHead}.
\end{enumerate}

\subsection*{c) Operaciones}


Costos: \texttt{addHead} $O(1)$ amortizado ($O(\texttt{size})$ solo cuando ejecuta resize); \texttt{removeTail} $O(1)$ worst-case; \texttt{get} $O(1)$.

\subsection*{d) \texttt{resize()}}
Al copiar en el orden lógico $i = 0, 1, \dots, \texttt{size}-1$ al inicio del nuevo arreglo, se preserva exactamente el mismo orden lógico, sin importar cómo estaban distribuidos físicamente antes. Tras el resize, $\texttt{tail} = 0$.

\subsection*{e) Invariantes}

\begin{enumerate}[label=\textbf{S\arabic*.}, leftmargin=1.2cm]
    \item $0 \le \texttt{size} \le \texttt{capacity}$ y $\texttt{capacity} \ge 1$.
    \item $0 \le \texttt{tail} < \texttt{capacity}$.
    \item \textbf{Mapeo consistente:} $\texttt{data}[\text{phys}(i)]$ es la posición lógica $i$ para todo $0 \le i < \texttt{size}$; en particular, $\texttt{data}[\texttt{tail}]$ es la cola y $\texttt{data}[\text{phys}(\texttt{size}-1)]$ es la cabeza.
    \item \textbf{Celdas libres:} las $\texttt{capacity} - \texttt{size}$ celdas no ocupadas por $\text{phys}(0..\texttt{size}-1)$ son \texttt{null}.
\end{enumerate}

Preservación:
\begin{itemize}[leftmargin=1.2cm]
    \item \textbf{\texttt{addHead}:} escribe en $\text{phys}(\texttt{size})$ (celda libre, S4) y aumenta \texttt{size} en 1; S1--S3 se mantienen.
    \item \textbf{\texttt{removeTail}:} lee $\texttt{data}[\texttt{tail}]$, lo pone a \texttt{null} y avanza \texttt{tail}; S1--S4 se mantienen.
    \item \textbf{\texttt{resize}:} reubica los elementos en orden lógico con $\texttt{tail} = 0$; S1--S4 se mantienen.
\end{itemize}

\subsection*{f) Traza}

Estado inicial: $\texttt{size} = \texttt{capacity} = 2$, cola $A$, cabeza $B$, $\texttt{tail} = 0$, $\texttt{data} = [A, B]$.

\renewcommand{\arraystretch}{1.3}
{\setlength{\tabcolsep}{4pt}
\begin{tabular}{@{}L{1cm}L{2cm}L{6cm}L{1cm}L{1cm}L{1.2cm}L{1.6cm}@{}}
\toprule
Mov. & Acción & \texttt{data} & \texttt{tail} & \texttt{size} & \texttt{cap.} & ¿Resize? \\
\midrule
--- & inicial &
$[A,B]$ &
0 & 2 & 2 & --- \\

$N_1$ & no come &
$[\text{null},B,C,\text{null}]$ &
1 & 2 & 4 & \textbf{Sí} \\

$N_2$ & come &
$[\text{null},B,C,D]$ &
1 & 3 & 4 & No \\

$N_3$ & no come &
$[E,\text{null},C,D]$ &
2 & 3 & 4 & No \\

$N_4$ & no come &
$[E,F,\text{null},D]$ &
3 & 3 & 4 & No \\

$N_5$ & come &
$[E,F,G,D]$ &
3 & 4 & 4 & No \\

$N_6$ & no come &
$[\text{null},E,F,G,H,\text{null},\text{null},\text{null}]$ &
1 & 4 & 8 & \textbf{Sí} \\

\bottomrule
\end{tabular}}
\renewcommand{\arraystretch}{1}



% Observaciones por movimiento (resize en N1, envoltura en N3--N4, resize con envoltura en N6).


%=====================================================================
\section*{P4 --- Análisis amortizado}

\subsection*{a) Aggregate Method}

La longitud inicial es $2$ y cada manzana suma $1$, pero como \texttt{addHead} va antes de \texttt{removeTail}, el tamaño pasa transitoriamente de $k$ a $k+1$. Por eso la longitud máxima transitoria es $n_{\max} = g + 3 = O(g)$. Con $C_0 = 2$ al ir duplicando la \texttt{capacity}, las capacidades son $2, 4, 8, \dots, 2^{r+1}$. Un resize se ejecuta cuando $\texttt{size} = \texttt{capacity}$, y en ese momento se copian exactamente $C_j$ elementos. Como la estructura nunca supera $n_{\max}$, hay $r = O(\log g)$ resizes y el total de copias es:
\[
\sum_{j=0}^{r-1} 2^{j+1} = 2^{r+1} - 2 < 2\,n_{\max} = O(g).
\]
\textbf{Resultado:} el total de copias es lineal en $g$, es decir $O(g)$.

\subsection*{b) Suma geométrica}

La capacidad solo cambia en un \texttt{resize}, y entre el \texttt{resize} $j$ y el $j+1$ vale $C_j$. El \texttt{resize} $j \to j+1$ se ejecuta cuando $\texttt{size} = C_j$, así que en ese instante el arreglo está lleno y se copian exactamente $C_j$ elementos. Por el proceso de duplicación, esos valores son $2, 4, 8, \dots, 2^{r}$. La suma es geométrica de razón $2$:
\[
\sum_{j=0}^{r-1} 2^{j+1} = 2^{r+1} - 2 < 2 \cdot 2^{r} = 2\,C_{r-1}.
\]
El último término domina a todos los anteriores, así que la suma es proporcional al último término y no a la cantidad de términos. Como $C_{r-1} = O(g)$, el total de copias es $O(g)$.

\subsection*{c) Costo total y amortizado de \texttt{addHead}}

El costo total real de $m$ movimientos se descompone en $O(m)$ por los \texttt{addHead} (escritura $+$ índices), $O(m)$ por los \texttt{removeTail} (leer, poner \texttt{null}, avanzar índice) y $O(g)$ por las copias de los \texttt{resizes}. Entonces $T(m) = O(m) + O(g) = O(m+g) = O(m)$, usando $g \le m$. El costo amortizado de \texttt{addHead} es $T(m)/m = O(m+g)/m = O(1)$. Aunque un \texttt{addHead} individual cueste $O(n)$ en el peor caso por un \texttt{resize}, su costo amortizado es constante.

\subsection*{d) Papel de $m$ y $g$}

$m$ y $g$ miden cosas distintas. $m$ controla el trabajo por operación: cada \texttt{MOVE} hace un \texttt{addHead} y, si no come, un \texttt{removeTail}, ambos $O(1)$; esta es la parte del costo que depende de cuántas operaciones se hacen. $g$ controla el crecimiento: la longitud máxima transitoria es $g+3$, la capacidad final es $O(g)$, hay $O(\log g)$ \texttt{resizes} y $O(g)$ copias; esta es la parte del costo que depende del tamaño alcanzado. La descomposición es $T(m) = O(m) + O(g)$. Como $g \le m$, el total es $O(m)$, pero las dos fuentes son conceptualmente independientes: con $g = 0$ no hay resizes; con $g \approx m$ hay $O(\log m)$ \texttt{resizes} pero el costo sigue siendo lineal en $m$.

\subsection*{e) Método del potencial}

Se elige $\Phi(\texttt{size}, C) = 2\cdot\texttt{size} - C$. Justo después de un \texttt{resize} (con $\texttt{size} = C_{\text{old}}$, $C_{\text{new}} = 2C_{\text{old}}$) el potencial vale $\Phi = 0$; justo antes, con $\texttt{size} = C$, vale $\Phi = C$, suficiente para pagar las $C$ copias. El costo amortizado es $\hat{c} = c + \Delta\Phi$.

\[
\begin{array}{c|c|c}
\text{Operación}&\Delta\Phi&\hat c\\ \hline
\texttt{addHead} \text{ sin resize}&2&1+2=3\\
\texttt{addHead} \text{ con resize}&2-C&(C+1)+(2-C)=3\\
\texttt{removeTail}&-2&1-2=-1
\end{array}
\]

La suficiencia del potencial frente a \texttt{removeTail} se justifica porque esta operación \textbf{no libera capacidad}, así que no ejecuta \texttt{resizes} ni copias; solo reduce el potencial en $2$, es decir, devuelve crédito sin consumirlo. Como el potencial solo se gasta en los resizes (donde $\Delta\Phi \approx -C$) y cada resize ocurre cuando $\Phi = C$, el crédito acumulado siempre alcanza. En general, $\sum_{i=1}^{t} \hat{c}_i = \sum_{i=1}^{t} c_i + \Phi_t - \Phi_0$, con $\Phi_0 = 2$ y $\Phi_t \le 2\cdot\texttt{size}_t = O(g)$, de donde $\sum c_i = O(m)$, coherente con el método agregado. El costo amortizado por operación es $O(1)$.


%=====================================================================
\section*{P5 --- The Adversarial Snake}

\subsection*{a) El adversario no cambia el número de copias}

Al no tener un adversario, un \texttt{resize} solo se ejecuta cuando $\texttt{size} = \texttt{capacity}$, y en ese instante se copian exactamente $C_j$ elementos. El adversario controla cuándo ocurre cada \texttt{resize}, pero no cuántos ni de qué tamaño (la secuencia de capacidades $C_0, C_1, C_2, \dots$ depende solo de la política y de $n$). Por tanto, el total de copias $\sum C_j$ es el mismo sin importar el momento del crecimiento.

Para aumentar las copias, el adversario necesitaría poder ejercutar un \texttt{resize} con $\texttt{size} < \texttt{capacity}$, es decir, controlar también la política de capacidad o contraer la estructura para volver a expandirla. Sin contracción y con política fija, el número de copias queda determinado por $n$ y $\alpha$.

\subsection*{b--e) Comparación de políticas}

\begin{itemize}[leftmargin=1.2cm]
    \item \textbf{$C+1$:} se copian $2+j$ en el resize $j$, luego $\sum_{j=0}^{n-3}(2+j) = \Theta(n^2)$.
    \item \textbf{$2C$:} serie geométrica de razón $2$, $\sum_{j=0}^{r-1} 2^{j+1} = 2^{r+1}-2 < 2n = O(n)$, y como la última copia es $\Theta(n)$, también $\Omega(n)$. Total $\Theta(n)$.
    \item \textbf{$\lceil 1.5C\rceil$:} serie geométrica de razón $1.5$; el último término domina, $\sum_{j=0}^{r-1} C_j = (C_r - C_0)/(1.5-1) = \Theta(n)$.
\end{itemize}

En los tres casos el costo amortizado por crecimiento es el total de copias dividido entre las $n$ inserciones que provocaron crecimiento. El espacio sin usar tras el resize es $C_{\text{new}} - (C_{\text{old}}+1)$, y su fracción se mide respecto a $C_{\text{new}}$.

\renewcommand{\arraystretch}{1.4}
{\setlength{\tabcolsep}{4pt}
\begin{tabular}{@{}L{1.7cm}L{2.1cm}L{2cm}L{2.4cm}L{2.8cm}@{}}
\toprule
Política & \# Resizes & Total copias & Costo amortizado & Espacio sin usar tras resize \\
\midrule
$C+1$              & $\Theta(n)$      & $\Theta(n^2)$ & $\Theta(n)$ & $0$ \ ($0\%$) \\
$2C$               & $\Theta(\log n)$ & $\Theta(n)$   & $\Theta(1)$ & $\approx C/2$ \ ($\approx 50\%$) \\
$\lceil 1.5C\rceil$ & $\Theta(\log n)$ & $\Theta(n)$   & $\Theta(1)$ & $\approx C/3$ \ ($\approx 33\%$) \\
\bottomrule
\end{tabular}}
\renewcommand{\arraystretch}{1}

\subsection*{f) Trade-off}

$C+1$ ahorra memoria ($0\%$ libre) pero es lento ($\Theta(n^2)$ copias, $\Theta(n)$ amortizado). $2C$ es rápido ($\Theta(1)$ amortizado) pero desperdicia $\approx 50\%$ de la capacidad. $\lceil 1.5C\rceil$ es el punto intermedio: mismo orden asintótico que $2C$ ($\Theta(\log n)$ resizes, $\Theta(n)$ copias, $\Theta(1)$ amortizado) pero con más resizes en constante, a cambio de desperdiciar solo $\approx 33\%$.

\subsection*{Familia $C_{\text{new}} = \lceil \alpha C \rceil$}

Para que la capacidad crezca se necesita $\alpha > 1$. La suma de capacidades es geométrica de razón $\alpha$; como $\alpha > 1$, el último término domina:
\[
\sum_{j=0}^{r-1} C_j = \frac{C_r - C_0}{\alpha - 1} = \Theta\!\left(\frac{n}{\alpha - 1}\right),
\qquad r = \Theta(\log_\alpha n).
\]
Para $\alpha$ constante y $\alpha > 1$, el total de copias es $\Theta(n/(\alpha-1)) = \Theta(n)$ y el costo amortizado es $\Theta(1/(\alpha-1)) = \Theta(1)$. Por tanto, duplicar no es necesario: cualquier $\alpha > 1$ constante funciona. Lo que no funciona es $\alpha = 1 + o(1)$ (crecimiento subgeométrico), porque entonces las copias dejan de ser lineales en $n$.

\end{document}
